\section{Panorama competitivo de 2026: la capacidad de los modelos ha entrado en una "competencia multidimensional"}

\subsection{El momento DeepSeek y la divergencia de estrategias de código abierto entre China y Estados Unidos}

La conversación llama al impacto de DeepSeek R1 a principios de 2025 un "evento de aceleración". Su importancia no radica solo en la posición en algún ranking, sino en demostrar que un modelo de pesos abiertos de alta calidad puede reajustar rápidamente las expectativas del mercado: los desarrolladores reevalúan su dependencia de las API, las empresas reevalúan las rutas de despliegue privado y los investigadores reevalúan las líneas base reproducibles.

Nathan menciona en particular que el equipo chino no logró un avance puntual, sino que surgieron múltiples avances en paralelo: Z.AI, Minimax, Moonshot, entre otros, entregaron modelos utilizables en la misma ventana de tiempo. Detrás de este fenómeno no hay "un algoritmo secreto", sino la capacidad de las organizaciones de ingeniería para avanzar en paralelo bajo una competencia de alta presión. El juicio de Sebastian también es directo: \textit{``ideas are not proprietary; resources are''}. En otras palabras, la diferencia suele estar en los recursos de ejecución, no en la exclusividad de los conceptos.

\begin{importantbox}{La licencia de código abierto se está convirtiendo en un arma comercial}
Antes, al hablar de código abierto, el foco solía estar en la "transparencia técnica"; ahora el foco más realista está en la "forma de entrar al mercado".
\begin{itemize}[leftmargin=1.5em]
    \item Los pesos abiertos con menos restricciones en la licencia pueden ampliar rápidamente el alcance entre los desarrolladores;
    \item para las empresas, un despliegue controlable importa más que el "mejor benchmark";
    \item para los proveedores de modelos, la apertura en sí misma puede canjearse por posicionamiento en el ecosistema y por palancas de comercialización posteriores.
\end{itemize}
\end{importantbox}

\begin{warningbox}{"Código abierto = almuerzo gratis" es un malentendido}
Los pesos abiertos reducen el costo de acceso, pero no eliminan el costo de ingeniería. Las empresas siguen teniendo que pagar los recursos de inferencia, la limpieza de datos, el sistema de evaluación, la gobernanza de registros, las auditorías de cumplimiento y el mantenimiento para revertir modelos. Lo verdaderamente difícil es convertir la capacidad del modelo en un servicio estable, no descargar el checkpoint.
\end{warningbox}

\subsection{Las ventajas reales de los principales modelos de código cerrado}

El juicio de la entrevista sobre el bando de código cerrado es mesurado. La ventaja de OpenAI se describe como "llevar continuamente nuevas direcciones de investigación a producción de forma rápida"; la ventaja de Google se describe como "cómputo e infraestructura integrados verticalmente"; la ventaja de Anthropic se describe como "haber formado un consenso de producto fuerte en el escenario de coding". Las tres son, en esencia, tipos distintos de capacidad organizacional.

\begin{knowledgebox}{Por qué en el mismo periodo surge un "desfase de percepción"}
El nivel de discusión en redes sociales y la escala real de usuarios a menudo se desfasan:
\begin{itemize}[leftmargin=1.5em]
    \item algunos modelos tienen mucho revuelo en X, pero eso no significa una alta penetración empresarial;
    \item algunos modelos tienen una reputación mediocre entre los desarrolladores, pero crecen muy rápido en los puntos de entrada masivos;
    \item lo que realmente determina la participación de mercado a largo plazo es el punto de entrada predeterminado, la estructura de precios, la experiencia de latencia y la estabilidad, no solo la posición en un ranking puntual.
\end{itemize}
\end{knowledgebox}

\subsection{Resumen de la sección}

La competencia de 2026 ya no es una guerra binaria de "código cerrado contra código abierto", sino un juego combinado en cinco dimensiones: licencia, costo, punto de entrada, herramientas y marca. El momento DeepSeek solo hizo visible y aceleró esta tendencia.

